\documentclass[11pt,letterpaper]{article}
\usepackage[margin=0.7in]{geometry}
\usepackage[T1]{fontenc}
\usepackage{lmodern}
\usepackage{amsmath,amssymb,booktabs,array,xcolor,listings,hyperref,fancyhdr}
\definecolor{accent}{RGB}{25,65,100}
\hypersetup{colorlinks=true,linkcolor=accent,urlcolor=accent}
\setlength{\parindent}{0pt}
\setlength{\parskip}{4pt}
\setlength{\headheight}{14pt}
\pagestyle{fancy}\fancyhf{}\fancyhead[L]{\small CSI2110 -- Study Notes}\fancyhead[R]{\small Fall 2026}\fancyfoot[C]{\thepage}
\lstset{basicstyle=\small\ttfamily,breaklines=true,columns=fullflexible,keepspaces=true,frame=single,rulecolor=\color{accent!25},backgroundcolor=\color{accent!3},showstringspaces=false,tabsize=2}
\newcommand{\sourcepages}[1]{\textit{Source: PDF pages #1.}}
\newcommand{\keyidea}[1]{\par\textbf{Key idea.} #1\par}
\newcommand{\cost}[1]{\ensuremath{O(#1)}}
\newcommand{\method}[1]{\texttt{#1}}

\title{\textbf{Lecture 0: Introduction}\\[3pt]\large CSI2110 -- Data Structures and Algorithms}
\author{Lecture study notes}\date{}\begin{document}\maketitle
\textbf{Source document:} \path{../lec00-introduction.pdf}.
These notes synthesize the source and add worked explanations. Page references use PDF page numbers, which can differ from printed slide numbers. Pseudocode is normalized for readability and boundary cases.
\tableofcontents\newpage

\section{What this course studies}
\sourcepages{2--13}
An \textbf{algorithm} describes steps for solving a problem; a \textbf{data structure} stores and organizes data so the required operations can run efficiently. CSI2110 focuses on choosing structures, implementing algorithms in Java, and explaining their time and space costs mathematically.

A phone book stores names and phone numbers and supports insert, delete, and lookup. An unordered list, ordered list, or tree can represent the same information while giving different operation costs. A map of buildings and distances naturally uses a graph and supports questions such as shortest routes.

\keyidea{Distinguish logical organization from physical implementation. A list, tree, graph, set, or map describes relationships; arrays, linked nodes, and adjacency lists describe ways to store those relationships.}

\section{Arrays}
\sourcepages{29--33}
An array is an indexed sequence with a fixed capacity. If the capacity is $N$, legal indices are $0,\ldots,N-1$. An array of primitives stores primitive values; an array of objects stores references. Java initializes an \method{int[]} with zeroes and an object-reference array with \method{null}.
\par\noindent\begin{minipage}{\linewidth}
\begin{lstlisting}[language=Java]
int[] x = {3, 1, 7, 5};
int[] y = new int[4];  // {0, 0, 0, 0}
\end{lstlisting}
\end{minipage}\par
\textbf{Size versus capacity:} $n$ is the number of logically stored entries; $N$ is the number of allocated slots. They need not be equal.

\subsection{Access and updates}
\begin{itemize}
\item Reading or replacing the value at a known index uses direct access: $O(1)$.
\item To insert at index $i$ in an ordered sequence, shift entries $i,\ldots,n-1$ one slot right. Shift from right to left so values are not overwritten.
\item To remove index $i$, shift entries $i+1,\ldots,n-1$ one slot left.
\item If an insertion would exceed capacity, allocate a larger array and copy existing entries. A resize takes $\Theta(n)$ time.
\end{itemize}
\textbf{Example:} starting with $[3,1,5]$, inserting $7$ at index $2$ produces $[3,1,7,5]$ after shifting $5$. Removing index $1$ then gives $[3,7,5]$.

\section{Singly linked lists}
\sourcepages{34--38,43}
Each node stores an element and a reference to the next node. The head refers to the first node; the last node has a null next reference. Nodes need not occupy consecutive memory locations.
\par\noindent\begin{minipage}{\linewidth}
\begin{lstlisting}
head -> [3 | next] -> [1 | next] -> [7 | next] -> [5 | null]
\end{lstlisting}
\end{minipage}\par
\par\noindent\begin{minipage}{\linewidth}
\begin{lstlisting}[language=Java]
class Node<E> {
    E element;
    Node<E> next;
    Node(E e, Node<E> n) { element = e; next = n; }
}
\end{lstlisting}
\end{minipage}\par
Generics allow the same structure to hold different reference types, such as \method{Integer} or \method{String}. Java generics use wrapper types rather than primitives.

\subsection{Local insertion and deletion}
Given a predecessor \method{pred}, insert immediately after it:
\par\noindent\begin{minipage}{\linewidth}
\begin{lstlisting}
newNode = new Node(e, pred.next)
pred.next = newNode
\end{lstlisting}
\end{minipage}\par
The order matters: save the successor before changing the predecessor's next reference.

Given the predecessor of a node to delete:
\par\noindent\begin{minipage}{\linewidth}
\begin{lstlisting}
pred.next = pred.next.next
\end{lstlisting}
\end{minipage}\par
Both changes take $O(1)$ once the predecessor is known. Finding the node at index $i$ requires traversal and takes $O(i+1)$, hence $O(n)$ in the worst case. Removing the head is a separate case: \method{head = head.next}.

\subsection{What a tail reference does}
A tail reference permits appending in $O(1)$: link the old tail to the new node and update the tail. Without a tail, appending requires traversal. A tail reference alone does \emph{not} permit constant-time removal of the tail, because its predecessor is not directly accessible. That operation still takes $O(n)$.

When the list becomes empty, both head and tail must be null. When inserting into an empty list, both must refer to the new node. A cached size field permits $O(1)$ \method{size()}.

\section{Doubly linked lists}
\sourcepages{39--42}
A doubly linked node stores an element, a previous reference, and a next reference. With references to both ends, traversal is possible in both directions.
\par\noindent\begin{minipage}{\linewidth}
\begin{lstlisting}
null <- [3] <-> [1] <-> [7] <-> [5] -> null
\end{lstlisting}
\end{minipage}\par
To insert $e$ between adjacent nodes \method{pred} and \method{succ}:
\par\noindent\begin{minipage}{\linewidth}
\begin{lstlisting}
n = new Node(pred, e, succ)
pred.next = n
succ.prev = n
\end{lstlisting}
\end{minipage}\par
To delete an interior node \method{x}:
\par\noindent\begin{minipage}{\linewidth}
\begin{lstlisting}
x.prev.next = x.next
x.next.prev = x.prev
\end{lstlisting}
\end{minipage}\par
Local updates take $O(1)$ when the relevant nodes are known. Index-based lookup still takes $O(n)$ in the worst case. Compared with a singly linked list, each node uses an extra reference.

\textbf{Sentinels:} a header before the first element and a trailer after the last element simplify updates. They store no user elements. An empty sentinel list has \method{header.next = trailer} and \method{trailer.prev = header}.

\section{Choosing an implementation}
\sourcepages{44}
\begin{center}\small
\begin{tabular}{@{}lccc@{}}\toprule
Operation & Array sequence & Singly linked$^*$ & Doubly linked$^*$\\\midrule
Access by index & $O(1)$ & $O(n)$ & $O(n)$\\
Insert/remove in middle & $O(n)$ & $O(1)$ local & $O(1)$ local\\
Append at end & $O(1)$ if room & $O(1)$ & $O(1)$\\
Remove last & $O(1)$ & $O(n)$ & $O(1)$\\
Storage & $O(N)$ & $O(n)$ & $O(n)$\\\bottomrule
\end{tabular}\end{center}
$^*$Assume a tail reference; singly linked local deletion requires the predecessor. An array resize costs $O(n)$. Lookup costs must be added to any local pointer change.

\section{Course logistics recorded in the slides}
\sourcepages{16--26}
This section records the supplied Fall 2026 slides; consult Brightspace for subsequent changes.
\begin{itemize}
\item Two lectures, one lab, and one tutorial per week. Lectures develop theory; labs practice Java; tutorials practice written reasoning.
\item Assessment weights: written assignments $15\%$, programming assignments $15\%$, midterm $30\%$, final $40\%$, plus up to $2\%$ lab bonus.
\item The slides list the midterm on November 7, 9:30--11:30 a.m.; programming assignments in October and November.
\item Let $A,P,M,F,B$ be weighted percentage-point contributions. If $(M+F)/70\geq0.5$, the stated grade is $A+P+M+F+B$; otherwise it is $100(M+F)/70$.
\item Assignments are individual and submitted through Brightspace. The stated late penalty is $30\%$ within 24 hours; submissions after 24 hours are not accepted.
\item Exams are closed-book. At least $50\%$ on the weighted exam component is needed to pass. An approved missed midterm shifts its weight to the final.
\item Labs are released Thursday at 4 p.m., due the following Thursday at 11:59 p.m., and require in-person TA grading within two weeks of the deadline. Session-specific details are on Brightspace.
\item Instructor office hours in the slides: Tuesday 3--4 p.m. and Thursday 2--3 p.m., SITE 4021. TA office hours are on Brightspace; email replies may take two business days.
\item The slides prohibit generative AI in submissions and require references for reused Internet code. These documents are study notes.
\end{itemize}

\section{Study checks}
\begin{enumerate}
\item Why is array access $O(1)$ but linked-list access by index $O(n)$? \textbf{Answer:} an index directly identifies an array slot; a list requires following links.
\item Is linked-list insertion always $O(1)$? \textbf{Answer:} only the pointer update is; locating an insertion position can take $O(n)$.
\item Why does a doubly linked list make tail deletion fast? \textbf{Answer:} \method{tail.prev} directly identifies the new tail.
\item What happens when an array is full? \textbf{Answer:} allocate more space and copy, or reject insertion in a fixed-capacity implementation.
\end{enumerate}

\end{document}
